% Generated draft: requires mathematical and source review.
\documentclass[11pt]{article}
\usepackage[margin=25mm]{geometry}
\usepackage{fontspec}
\setmainfont{Latin Modern Roman}
\usepackage{amsmath,amssymb,mathtools,unicode-math}
\setmathfont{Latin Modern Math}
\usepackage{enumitem,needspace,xurl,hyperref}
\setlength{\parindent}{0pt}
\setlength{\parskip}{5pt}
\setlength{\emergencystretch}{3em}
\newcommand{\N}{\mathbb{N}}
\newcommand{\Z}{\mathbb{Z}}
\newcommand{\Q}{\mathbb{Q}}
\newcommand{\R}{\mathbb{R}}
\newcommand{\C}{\mathbb{C}}
\newcommand{\F}{\mathbb{F}}
\newcommand{\A}{\mathbb{A}}
\newcommand{\PP}{\mathbb{P}}
\newcommand{\E}{\mathbb{E}}
\newcommand{\eps}{\varepsilon}
\newcommand{\dd}{\,\mathrm{d}}
\begin{document}
\section*{UnsolvedMath ID: 30003688}
\textbf{Source problem code:} catalogue entry
\subsection{Definitions and mathematical statement}
The supplied formulations do not determine an unambiguous mathematical problem. The following gives a conditional reading of the cleaned statement using customary numbering conventions; the source record does not specify these conventions.

Let \(n\in\Z\) with \(n\geq2\). Under this reading, \(\Sigma_n^0\) denotes the \(n\)-th existential level of the arithmetical hierarchy for subsets of \(\N\). For a nonempty family \(\mathcal S\subseteq\mathcal P(\N)\), an admissible numbering is a total surjection \(\nu:\N\to\mathcal S\) such that
\[
\{(e,x)\in\N^2:x\in\nu(e)\}\in\Sigma_n^0,
\]
using a fixed computable coding of pairs. Write \(\mathcal N_n(\mathcal S)\) for these numberings. For \(\mu,\nu\in\mathcal N_n(\mathcal S)\), define
\[
\mu\leq\nu
\quad\Longleftrightarrow\quad
\exists f:\N\to\N\text{ total computable}\ \forall e\in\N\;
\mu(e)=\nu(f(e)),
\]
and define \(\mu\equiv\nu\) by mutual reducibility. The corresponding Rogers semilattice \(\mathcal R_n(\mathcal S)\) is the quotient by \(\equiv\), ordered by the induced reducibility relation. A Friedberg numbering, under the usual convention, is an injective admissible numbering.

With these conditional definitions, the cleaned question asks whether
\[
\begin{gathered}
\forall n\in\Z,\ n\geq2,\quad
\exists\mathcal S\subseteq\Sigma_n^0\;
\exists a,b\in\mathcal N_n(\mathcal S):\\
0<|\mathcal S|<\infty,\qquad
\mathcal R_n(\mathcal S)=\{[a],[b]\},\qquad [a]\neq[b],\\
a\text{ is injective},\qquad
\forall\gamma\in\mathcal N_n(\mathcal S)\;
\bigl(\gamma\not\equiv a\Longrightarrow b\leq\gamma\bigr).
\end{gathered}
\]
Here \(\alpha=[a]\) and \(\beta=[b]\) distinguish semilattice elements from their numbering representatives.

\emph{Unresolved qualification.} A total injective numbering with domain \(\N\) cannot have finite range. Thus the supplied finite-family requirement and its Friedberg requirement are incompatible under the conventions just described. The record does not establish whether a different numbering domain or meaning of Friedberg is intended, or whether a formulation has been corrupted. This conditional display must therefore not be treated as a verified reconstruction of the intended problem.

The original statement additionally asks whether any family, without the finiteness restriction, has the referenced properties. Those properties are not specified there. It is consequently unresolved whether this second existence question concerns precisely the properties in the cleaned statement, and whether the cleaned universal quantifier over \(n\geq2\) applies to it.
\subsection{Short English statement}
For each integer \(n\geq2\), the cleaned record asks for a finite family of \(\Sigma_n^0\) sets with exactly two numbering degrees, one represented by a Friedberg numbering and the other below every numbering outside the first degree. The usual meaning of Friedberg numbering conflicts with finiteness. The original record also asks about unrestricted families but leaves the relevant properties unspecified. These issues require human mathematical review.
\subsection{Sources}
\begin{itemize}
\item Local registry: the supplied complete source record for permanent UnsolvedMath ID \(30003688\), titled \emph{Finite Families with Two-Element Rogers Semilattices}; exact source problem code \texttt{catalogue entry}, in the Oberwolfach Reports Open Problems collection. No individual proposer is supplied; the proposed year is recorded as 2018.
\item Supplied report citation: \emph{Computability Theory (2019)}, Oberwolfach Reports, DOI \texttt{10.4171/owr/2018/1}, pp.\ 5--41. Supplied source URL and literature-source link, labelled ``Oberwolfach report'': \url{https://doi.org/10.4171/owr/2018/1}. This report attribution is supplied evidence; its contents have not been independently inspected.
\item Literature assessment dated exactly \texttt{2026-08-21}: supplied dated curation, not independently verified. It labels the problem open and explains that the extracted record omits the properties and family being quantified, preventing verification of a definite claim; the open label was assigned by normalization. The background describes this as literature triage derived from supplied Oberwolfach curation and requests independent verification of time-sensitive status and cited claims.
\item The registry separately supplies the statement-status label \texttt{corrected\_verified} and reports manual semantic, typographic, and LaTeX review asserting that the cleaned statement is self-contained and adds no unsupported content. This is a supplied review claim, not independent verification, and does not resolve the formulation issues recorded here.
\end{itemize}
\end{document}
